\section{执行程序}
\begin{enumerate}
  \item 最复杂的一章，内容非常多，并且其中掺杂着不少的难点。
  \item \zd{4 个问题。启动、阻却、救济、特殊措施}。
  \item 申请执行的时效为 2 年，可以中止、中断。必须以\zd{生效法律文书}为根据（\zd{审判和执行没啥关系}）。\al{根据当事人执行和解协议制作调解书 ×}。
  \item 管辖。
        \begin{enumerate}
          \item 诉讼执行 \yjt 一审法院 + 被执行财产所在地法院。
          \item 非诉执行 \yjt 作出法院 + 被执行财产所在地法院。
          \item 其他执行 \yjt 被执行人住所地 + 被执行财产所在地法院。
          \item 执行管辖异议。
                \begin{enumerate}
                  \item 后立案的，撤销案件。
                  \item 收到执行通知之日起的 10 日内提。
                  \item 异议成立，撤销案件。
                  \item 对裁定不服的，上级复议。
                \end{enumerate}
        \end{enumerate}
  \item 启动方式。
        \begin{enumerate}
          \item 申请执行。\zd{申请人是生效法律文书确定的权利人 \hz 其继承人、权利承受人}。
          \item 移送执行。赡养/抚养/抚育费、刑事被害人、环保/检察院提起的公益诉讼。
        \end{enumerate}
  \item 通知。10 天内；可以委托当地法院代为执行。
  \item **********阻却**********
  \item 暂缓执行。
        \begin{enumerate}
          \item \zd{执行担保}。\al{甲申请暂缓执行，由乙提供担保，期限届满后甲仍不履行义务}。
                \begin{enumerate}
                  \item 开始计算暂缓执行期间、执行担保期间。
                  \item 暂缓期间转移、隐匿、变卖、毁损担保财产 \hz 暂缓期间届满。
                        \begin{enumerate}
                          \item \zd{申请}恢复执行 \yjt 直接执行担保财产 \hz 直接裁定执行担保人的财产。
                          \item 不考虑先诉抗辩权，不需要把担保人变更追加为被执行人。
                          \item 执行担保人的，以担保范围为限，以担保期间为限。
                        \end{enumerate}
                \end{enumerate}
        \end{enumerate}
  \item 执行和解。书面、口头达成均可。
        \begin{enumerate}
          \item 协议内容记入执行笔录，双方签字盖章 \yjt 裁定中止执行（等履行）。\al{法院调解，让乙接受甲玉石抵债 ×}。
                \begin{enumerate}
                  \item 履行完毕 \yjt 执行终结。\zd{瑕疵另诉}。\al{甲乙达成和解协议，以玉石抵债，当即交付。\xhx{后}乙发现是赝品}。
                  \item 拒不履行 \yjt 申请恢复执行 \hz \zd{就和解协议起诉}。\al{申请对和解协议强制执行 ×}。
                \end{enumerate}
          \item \zd{执行和解协议可诉}。起诉\zd{确认和解协议无效 \hz 撤销和解协议}。解决实践执行难问题，特殊规定。
          \item 产生的 3 种诉讼\zd{均由执行法院管辖}。
        \end{enumerate}
  \item 中止和终结。了解即可。执行中权利人\zd{撤回}申请后，法院应当裁定执行\zd{终结}。
  \item 执行回转。\zd{执行错}了。据以执行的法律文书被依法撤销、变更。\zd{裁定}要求\zd{返还}财产，拒不返还的强制执行。
  \item **********救济**********
  \item 异议。
        \begin{enumerate}
          \item 先做一对区分。上级复议 \yjt 法院来找我麻烦。同级复议 \yjt 我去找法院帮忙。
          \item 对执行\zd{行为的异议}。只能向上一级法院申请复议一次。\al{甲认为法院扣押钢琴（父亲的遗物）行为错误}。
          \item 案外人对执行\zd{标的的异议}。找执行法院 \yjt 如果是原法院下级 \yjt 找原法院 \yjt\, 15 日内审查。% https://q.uiver.app/#q=WzAsMTAsWzAsMSwi6KOB5a6a5Lit5q2iIl0sWzEsNCwi5omn6KGM5byC6K6u5LmL6K+JXFxcXO+8iOehruiupOS5i+ivie+8iSJdLFsxLDAsIuahiOWkluS6uuWvueaJp+ihjOagh+eahOeahOW8guiuriJdLFsyLDEsIumps+WbniJdLFswLDYsIueUs+ivt+S6uiJdLFszLDYsIuesrOS4ieS6uiJdLFsxLDIsIuWIpOaJpyJdLFszLDIsIuWGjeWuoSJdLFszLDAsIuS4ieaSpOeri+azleWkquaZmlxcXFzlj6rog73lho3lrqEiXSxbMyw0LCLooqvmiafooYzkuroiXSxbMiwwLCLmiJDnq4siLDFdLFsyLDMsIuS4jeaIkOeriyIsMV0sWzAsNl0sWzMsNl0sWzYsNywi5LiA6Ie077yI5Y6f5Yik5pyJ6K+v77yJIiwxXSxbNiwxLCLkuI3kuIDoh7TvvIjkuI7ljp/liKTml6DlhbPvvIkiLDFdLFs3LDhdLFs5LDEsIuS4jeWPjeWvue+8jOaXoOeLrOS4iSIsMl0sWzQsNSwi6K645Y+v5omn6KGM5LmL6K+JL+eUs+ivt+S6uuaJp+ihjOW8guiuruS5i+ivie+8iOa2iO+8iSIsMCx7Im9mZnNldCI6LTF9XSxbNSw0LCLmoYjlpJbkurrlvILorq7kuYvor4nvvIjnp6/vvIkiLDAseyJvZmZzZXQiOi0xfV0sWzEsNV0sWzEsNF0sWzksMSwi5Y+N5a+577yM5YWx5ZCM6KKr5ZGKIl1d
\[
  \begin{tikzcd}
    & \txt{案外人对执行标的的异议} && \txt{三撤立法太晚\\只能再审} \\
    \txt{裁定中止} && \txt{驳回} \\
    & \txt{判执} && \txt{再审} \\
    \\
    & \txt{执行异议之诉\\（确认之诉）} && \txt{被执行人} \\
    \\
    \txt{申请人} &&& \txt{第三人}
    \arrow["\txt{成立}"{description}, from=1-2, to=2-1]
    \arrow["\txt{不成立}"{description}, from=1-2, to=2-3]
    \arrow[from=2-1, to=3-2]
    \arrow[from=2-3, to=3-2]
    \arrow["\txt{一致（原判有误）}"{description}, from=3-2, to=3-4]
    \arrow["\txt{不一致（与原判无关）}"{description}, from=3-2, to=5-2]
    \arrow[from=3-4, to=1-4]
    \arrow["\txt{不反对，无独三}"', from=5-4, to=5-2]
    \arrow["\txt{许可执行之诉/申请人执行异议之诉（消）}", shift left=1, from=7-1, to=7-4]
    \arrow["\txt{案外人异议之诉（积）}", shift left=1, from=7-4, to=7-1]
    \arrow[from=5-2, to=7-4]
    \arrow[from=5-2, to=7-1]
    \arrow["\txt{反对，共同被告}", from=5-4, to=5-2]
  \end{tikzcd}
\]

                \begin{enumerate}
                  \item \yi 成立 \yjt 裁定中止。\er 不成立 \yjt 驳回。
                  \item \yjt\,\yi 判执一致（原判有误）\yjt 再审\footnote{理论来讲，这里可以去提三撤，但三撤立法太晚，只能再审。}；
                        \er 判执不一致（\zd{与原判无关}）\yjt \zd{执行异议之诉}（确认之诉）：
                        \begin{enumerate}
                          \item 核心争议。执行标的物到底是不是被执行人的；申请人和第三人在对抗。
                          \item 申请人 \yjt 第三人为被告 \yjt 提许可执行之诉/申请人执行异议之诉 - \zd{消}。为了继续执行。
                          \item 第三人 \yjt 申请人为被告 \yjt 提案外人异议之诉 - \zd{积}。为了阻碍执行。
                          \item 被执行人不反对/反对 \yjt 无独三/共同被告。执行标的物最后是不是他的无所谓了。\zd{看热闹}。
                          \item 15 日；书面；执行法院一审普通程序；\zd{谁阻碍执行谁承担举证责任}；\al{查扣冻 √ 处分 ×}。
                        \end{enumerate}
                \end{enumerate}
        \end{enumerate}
  \item *********特殊措施。主体特殊 + 标的特殊**********
  \item 特殊主体的执行。
        \begin{enumerate}
          \item 执行中变更、追加当事人。
                \al{执行中甲死亡，继承人乙表示继承，其他继承人均表示放弃继承 \yjt 变更乙为被执行人}。
                \begin{enumerate}
                  \item 追加。到底追加变更谁为被执行人 \yjt 哪些人对被执行债务要\zd{承担责任}的。民商结合。
                  \item 救济。只有“身份” \yjt 上级复议（\zd{不解决实体争议}）；身份 + 涉及\zd{实体纠纷} \yjt 另行起诉。
                        \al{甲设立一人有限公司，乙申请执行时发现公司账户有 100W，且发现财产混同，遂申请追加甲的公司为被执行人。
                          甲如果对法院裁定不服, 可以另行起诉 √，可以向上一级法院申请复议 ×}。
                \end{enumerate}
          \item 参与分配。\zd{多个债权人}对一个债务人。除法人之外的自然人和其他组织。
                \al{甲申请执行张三的财产，乙丙丁申请参与分配。甲乙不服分配方案提出异议，法院对分配方案进行修订，丙丁均反对}。
                \begin{enumerate}
                  \item 分配方案。债权人、被执行人有权提出异议 \yjt 通知没提异议的人 \yjt\, ... \yjt 修订方案 \yjt 有人反对。
                  \item 异议者告反对者 \yjt \zd{分配方案异议之诉}。\al{甲乙以丙丁为被告向执行法院提起诉讼}。
                  \item 不存在上诉、复议等问题。
                \end{enumerate}
          \item 执行转破产。不成功 \yjt 恢复执行 \yjt 查扣冻，\zd{先到先得}。
        \end{enumerate}
  \item 特殊标的的执行。\zd{即便某个做法很合理，法院也不能自己直接作出！！！！}
        \begin{enumerate}
          \item 代位申请执行债权。\zd{对被执行人到期债权的执行}。有点类似督促程序。
                \begin{enumerate}
                  \item \zd{申请}法院向次债务人发出\zd{通知}，15 天内书面 \hz 口头提出\zd{实质异议}（\al{还了；过时效}）。
                        \begin{enumerate}
                          \item 不还钱 \hz 不提异议 \yjt 强制执行；罚款、拘留；在无法\zd{清偿范围内}承担\zd{连带}责任。
                          \item 提了 \yjt 不执行。\zd{法院不对异议审查}。\al{法院应当对提出的异议进行审查 ×}。
                        \end{enumerate}
                  \item 只能代位 1 次，不允许再代位。\al{可就次债务人对他人享有的到期债权强制执行 ×}。
                \end{enumerate}
          \item \zd{共有物}。\al{法院查扣了甲的汽车，甲的两个哥哥主张该汽车三兄弟共有。法院经常属实}。
                \begin{enumerate}
                  \item 可以查扣冻（有被执行人份额）\yjt 及时通知共有人。
                  \item 析产分割 \yjt 属于实体纠纷 \yjt\, \yi 协议分割 + 申请人同意；\er （被执行人/申请人代位）提析产诉讼 \yjt 执行中止。
                \end{enumerate}
          \item \zd{特定物}。执行不能。每年都考，越来越难。
                \begin{enumerate}
                  \item \zd{必须要执行原物}，否则执行无意义。
                  \item \zd{灭失毁损}（扩大解释为\zd{给付不能}）\yjt 折价赔偿 \yjt 属于实体争议 \yjt 协议 \hz 起诉 \yjt 执行终结。
                  \item \al{可以申请法院执行 xx 价值 xx 元的其他财产 ×}。
                \end{enumerate}
          \item \zd{保留所有权}买卖合同中的标的物。
                \begin{enumerate}
                  \item 可以查扣冻。
                  \item 被执行人将其财产\zd{出卖}给第三人。考试考了很多次。
                        \al{甲卖画给乙，约定乙首付 10W 后交付字画，付清前保留所有权。后丙因借款纠纷起诉甲，法院执行中扣押了甲的画}。
                        \begin{enumerate}
                          \item 第三人\zd{向法院}给付剩余价款（\al{提供担保 ×}） \yjt 法院解除查扣冻 \yjt 继续\zd{履行合同}。实质为第三人代为履行。
                        \end{enumerate}
                  \item 被执行人\zd{购买}第三人财产。没考过。
                        \al{甲卖画给乙，约定乙首付 10W 后交付字画，付清前保留所有权。后丙因借款纠纷起诉乙，法院执行中扣押了甲的画}
                        \begin{enumerate}
                          \item 要钱。优先受偿（登记对抗）。
                          \item 要物。所有权人如要主张物的权利，应以案外人身份提\zd{案外人对执行标的的异议}。
                        \end{enumerate}
                \end{enumerate}
          \item 赔礼道歉、恢复名誉。判决的\zd{主文}在媒体上\zd{公告}。
          \item 逾期/迟延履行。
                \begin{enumerate}
                  \item \zd{行为义务}。迟延履行\zd{金}；\zd{双倍}补偿损失。
                  \item \zd{金钱义务}。\zd{加倍}支付迟延履行期间\zd{债务利息}。
                \end{enumerate}
        \end{enumerate}
  \item 强制执行措施。
        \begin{enumerate}
          \item 财产报告制度。了解即可。
          \item 限制出境；限制高消费。
          \item 记录到征信系统。
          \item 媒体公布不得强制执行的财产。
        \end{enumerate}
\end{enumerate}
